%% Supplementary: cohort and arm details moved from Section 5 of the main text.
\section{Additional details of the cohorts and arms}
\label{sup:setup}

\subsection{Cohorts}
\label{sup:cohorts}

\paragraph{Camelyon17-WILDS} The cohort consists of $96\times96$~px
lymph-node patches labeled tumor or normal from five hospitals
\citep{koh2021wilds,bandi2019camelyon17}. In fold $t$, the test split is a
class-balanced sample of 25,000 patches of hospital $t$, the OOD-validation
split 12,000 patches of hospital $(t-1)\bmod5$, and the training split 30,000
class-balanced patches of the other three hospitals. ID validation (12,000
patches) holds out $\max(1,\lfloor n_c/5\rfloor)$ of the $n_c$ patients of
every training hospital $c$. The patients are redrawn with a seeded generator
until the held-out set has at least 1,000 patches of each class and at most
40\% of the hospital's tumor patches, because patients contribute between 5
and 56,666 tumor patches. The 64~px cache is the central crop of the native
patch.

\paragraph{Canine cutaneous squamous cell carcinoma} Each of 44 slides was
digitized with five scanners \citep{wilm2023canine}: cs2 (Aperio ScanScope
CS2), gt450 (Aperio GT450), p1000 (3DHistech Pannoramic~1000), nz20 (Hamamatsu
NanoZoomer 2.0-HT) and nz210 (NanoZoomer S210). A domain is a scanner. The
slides are randomly split into five groups. Fold $t$ tests scanner $t$ on slide
group $t$, validates on scanner $(t-1)\bmod5$ with slide group $(t-1)\bmod5$,
and trains on the remaining three scanners and slide groups. A test domain is
therefore a new scanner and a set of slides that is absent from training under
every scanner. Four training slides (15\%) form the ID-validation split. Per
image, about 420 candidate centers are drawn, half in carcinoma polygons and
half in non-tumor tissue polygons, allocated in proportion to polygon area. Negative centers that fall inside any tumor polygon
are dropped. A patch is kept if it lies inside the image at both sizes and its
64~px crop passes an intensity filter (mean $\le225$, SD $\ge10$), and
duplicate centers are removed. Magnification is normalized to cs2: the factor
of each scanner is the median over slides of the square root of its
annotation-area ratio to cs2. Crops of 64/65/71/73/62~px (cs2, p1000, nz20,
nz210, gt450) are resized bilinearly to 64~px, and crops of
96/98/107/110/92~px to 96~px. Factors recomputed without each fold's test and
validation slides give the same crop sizes.

\paragraph{MIDOG 2021} The training set of the MIDOG~2021 challenge
\citep{aubreville2023midog} contains 150 annotated human breast cancer images
from one laboratory, 50 from each of three scanners (Hamamatsu NanoZoomer XR,
NanoZoomer S360 and Aperio ScanScope CS2). A domain is a scanner ($K=3$), and
there is no OOD-validation domain. Patches are centered on the annotated
bounding boxes of mitotic figures and non-mitotic look-alikes. ID validation
holds out 15\% of the images of every training scanner (eight per scanner),
fixed per scanner. The TIFF resolutions are 0.2263--0.2274 (XR), 0.2298 (S360)
and 0.2533 (CS2)~\textmu m/px, a difference of about 11\%. The primary cache
therefore normalizes every image to 0.25~\textmu m/px: it crops
$\mathrm{round}(P\cdot0.25/\mathrm{mpp})$~px (63--71 and 95--106~px) and
resizes bilinearly to $P\in\{64,96\}$. The un-normalized cache keeps the
magnification difference as part of the shift, as in the challenge. It is run
at the compact tier as a sensitivity analysis.

\paragraph{MIDOG++} MIDOG++ \citep{aubreville2023midogpp} provides 503
annotated 2~mm$^2$ regions of interest from seven tumor types in two species
(Table~\ref{tab:midogpp}). Its 150 breast cancer images have pixel data
identical to the MIDOG~2021 images. A domain is a tumor type ($K=7$), and
laboratory, scanner and species vary with it. Task, patch centering, scale
normalization and ID validation follow MIDOG~2021. For the 3DHistech images,
whose $x$ and $y$ resolutions differ by 0.4\%, the geometric mean is used.

\begin{table}[width=\linewidth,pos=t]
\caption{MIDOG++ domains \citep[scanners and laboratories from Table~1
of][]{aubreville2023midogpp}. Counts are the annotated mitotic figures and
non-mitotic look-alikes used in this study; resolution is from the TIFF tags.
Fold $t$ tests domain $t$ and validates on domain $(t-1)\bmod 7$.}
\label{tab:midogpp}
\footnotesize
\begin{tabular*}{\tblwidth}{@{\extracolsep{\fill}}llrrr>{\raggedright\arraybackslash}p{2.75cm}>{\raggedright\arraybackslash}p{2.2cm}l@{}}
\toprule
Tumor type & Species & Images & Mitotic & Look-alike & Scanner & Laboratory & \textmu m/px \\
\midrule
Mast cell tumor & canine & 50 & 2,327 & 1,366 & Aperio ScanScope CS2 & FU Berlin & 0.2533 \\
Lung cancer & canine & 44 & 855 & 951 & 3DHistech Pannoramic Scan II & VMU Vienna & 0.2487 \\
Lymphosarcoma & canine & 55 & 3,959 & 4,257 & 3DHistech Pannoramic Scan II & VMU Vienna & 0.2487 \\
Soft tissue sarcoma & canine & 100 & 1,286 & 2,375 & 3DHistech Pannoramic Scan II & AMC New York, VMU Vienna & 0.2487 \\
Breast cancer & human & 150 & 1,721 & 2,714 & Hamamatsu XR, S360; Aperio CS2 & UMC Utrecht & 0.2263--0.2533 \\
Melanoma & human & 49 & 1,150 & 925 & Hamamatsu XR & UMC Utrecht & 0.2263--0.2274 \\
Neuroendocrine tumor & human & 55 & 639 & 1,761 & Hamamatsu XR & UMC Utrecht & 0.2263--0.2274 \\
\midrule
Total & & 503 & 11,937 & 14,349 & & & \\
\bottomrule
\end{tabular*}
\end{table}

\paragraph{Mirror padding} MIDOG objects closer to the image border than
half a crop are cut from a reflect-padded copy of the image. The object
therefore stays centered, and its 64 and 96~px patches are concentric. At
96~px this affects 111 of 4,435 MIDOG~2021 patches and 617 of 26,286 MIDOG++
patches, at similar rates in both classes (MIDOG++: 2.43\% of positives,
2.28\% of negatives). The padding of each row is stored for a sensitivity
analysis that excludes these patches.

\paragraph{Unlabeled pools} The pool used for self-supervised pre-training
comes only from the training domains and excludes ID-validation groups, the
validation domain and the test domain. The arms that use it therefore remain
domain generalization methods and do not adapt to the test domain. On
Camelyon17 the pool consists of 40,000 patches of non-ID-validation training
patients that are not in the training split. On the other cohorts it consists
of random tissue crops (mean $<220$, SD $>12$) from non-ID-validation training
images. The crops are drawn with image-specific seeds, at most 800 per canine
image and 300 per MIDOG image, and subsampled to 40,000. MIDOG~2021 yields
only 25,200.


\subsection{Self-supervision controls and stain normalization}
\label{sup:arms}

\paragraph{Controls for self-supervision} The augmentation-only control
applies SimCLR's view augmentation to labeled batches without the contrastive
loss. It separates the effect of the augmentations from that of pre-training.
As in the SimCLR views, crop and flip parameters are drawn per batch and color
factors per image. The compute-matched control gives ERM as many images through
the encoder as SimCLR receives in pre-training and fine-tuning together.
Pre-training runs 6 epochs of $\lfloor|U|/256\rfloor$ steps on two views of 256
images, so $I_{\text{SSL}}=479{,}232$ for $|U|=40{,}000$ and 301,056 for
MIDOG~2021. SimCLR is pre-trained separately for every fold and seed, with SGD
(momentum 0.9, weight decay $10^{-4}$), a cosine schedule and projection heads
of 256--256--128 (C) or 2048--2048--128 (S). It is then fine-tuned with the ERM
recipe.

\paragraph{Stain normalization} As in common practice, Macenko's stain matrix
\citep{macenko2009} is estimated once per slide. A group is a slide imaged on
one scanner: a Camelyon17 slide, a MIDOG image, a canine slide on one of its
three scanners, or a PatchCamelyon whole-slide image. For each group, up to 256
tiles are drawn in a content-hashed order that does not depend on row order.
Their tissue pixels (optical density $\ge0.15$ in every channel; at most
200,000, subsampled with a fixed seed) are pooled. Macenko's estimator is then
applied once: the plane of the two leading eigenvectors, then the 1st and 99th
angular percentiles. The two directions are labeled hematoxylin and eosin by
matching them to the reference. Concentrations are the least-squares solution
clamped at zero, because without the clamp eosin-rich pixels received large
negative hematoxylin values on pale tiles. The group's maximal concentrations
are the per-stain 99th percentiles over its pooled pixels. Every tile is
rescaled to the reference maxima and reconstructed with the reference matrix. A
group with too little tissue keeps the reference. The reference is fitted per
fold with the same estimator on every training group: its stain matrix is the
elementwise median of the aligned group matrices, and its maxima are the median
group maxima. No labels are used. The method is transductive, however: each
test slide is normalized from its own tiles, which are the benchmark's patches
from annotated regions rather than a label-independent sample of the slide.
Section~\ref{app:macenko} reports how much the estimate depends on that tile
set, and a per-tile variant that estimates the stain matrix on every tile
separately. Our H\&E jitter perturbs RGB color; it is not an augmentation in
HED space.

